\chapter{Making sensors comparable: band filtering, interpolation and normalisation}\label{ch:bandpipeline}

\begin{onething}
Eight stages turn a sensor's own bands into $165$ bands that mean the same thing on every
hyperspectral sensor, using nothing more advanced than the equation of a straight line (and one
shape-preserving curve). Inside the model, a second and different ``normalisation'' rescales each
band to a z-score, and its exact inverse turns the model's output back into reflectance.
\end{onething}

\section{The problem}\label{sec:bandpipeline:problem}

\Cref{ch:sensors} fixed the destination: $165$ wavelengths from $460$ to $2{,}450$ nm. This chapter
is the journey. A PRISMA cube arrives with $239$ bands at PRISMA's wavelengths, some flagged bad by
the sensor, some inside atmospheric windows, some noisy at detector edges, and pixels that are
valid in some bands and not others. To reach the grid you must decide which bands to trust, decide
which pixels to keep, fill the holes in the kept pixels' spectra, and then estimate each spectrum's
value at the grid wavelengths.

The obvious shortcut is to skip the first three decisions and interpolate the raw spectrum straight
onto the grid. It fails, because an interpolator trusts every number it is given: a bad band
becomes a spike in the resampled spectrum, an invalid zero drags its neighbours down, and a
detector-edge artefact is spread across two grid bands. Interpolation is only as good as its
inputs, so the pipeline cleans first and interpolates last.

The pipeline is driven entirely by one configuration object, \code{BandFilterConfig}
(\file{app/models/dataset/vendables.py}, lines 71--138), and run by each sensor's dataset builder;
the shared composition is visible in one function, \code{to_common_grid}
(\file{app/utils/data_transformations/common_grid_pipeline.py}).

\begin{table}[htbp]\centering
\caption{The eight stages and their defaults. Stages 1--4 are \texttt{SpectralBandFilter}
(\texttt{spectral\_band\_filter.py}); stage 7 is \texttt{interpolate\_spectral\_gaps}
(\texttt{spectral\_interpolator.py}); stage 8 is \texttt{resample\_to\_common\_grid}
(\texttt{spectral\_resampler.py}).}\label{tab:bandpipeline:stages}
\small
\begin{tabular}{@{}clp{0.33\linewidth}l@{}}\toprule
Stage & Acts on & What it does & Setting (default)\\\midrule
1 & bands & drop bands the sensor flagged bad & always\\
2 & bands & drop bands inside atmospheric windows & \texttt{exclusion\_ranges} (5 windows)\\
3 & bands & drop the first and last bands of each detector & \texttt{edge\_bands\_to\_trim} ($3$)\\
4 & bands & drop bands with too few valid pixels & \texttt{min\_valid\_pixel\_pct} ($20.0$)\\
5 & pixels & invalidate pixels flagged by quality masks (EnMAP only) & \texttt{quality\_masks\_to\_apply}\\
6 & pixels & invalidate pixels with too many invalid bands & \texttt{max\_invalid\_voxel\_fraction} ($0.4$)\\
7 & values & fill the remaining holes along wavelength (PCHIP) & --\\
8 & values & resample onto the common grid (linear) & \texttt{common\_wavelength\_grid}\\\bottomrule
\end{tabular}
\end{table}

After the eight stages, the \code{band_filter_apply} Action applies one more step, the nearest-valid
pixel fill (\cref{sec:bandpipeline:fill}), and pickles the result.

\begin{sourcenote}[: ``four stages'' or ``eight'']
The docstring of \code{BandFilterConfig} says it ``controls four filtering stages''; the courses and
\file{docs/03-data-pipeline.md} speak of eight. Both are right about different things: the four are
the band-dropping stages inside \code{SpectralBandFilter}, and the eight count the pixel and value
stages the builders run afterwards with the same configuration. This book uses the eight.
\end{sourcenote}

\section{Stages 1--4: deciding which bands to trust}\label{sec:bandpipeline:bands}

\code{SpectralBandFilter.get_good_band_indices} walks the sensor's bands once and keeps band $k$
only if it passes all four tests, in order:
\begin{lstlisting}
for i in range(len(self._wavelengths)):
    if self._flags[i] != 1: continue                                   # 1: sensor flag
    wl = self._wavelengths[i]
    if any(lo <= wl <= hi for lo, hi in self._exclusion_ranges): continue  # 2: window
    if i in edge_set: continue                                         # 3: detector edge
    if self._validity_scores[i] < self._min_valid_pixel_pct: continue  # 4: coverage
    good.append(i)
\end{lstlisting}

\textbf{Stage 1, the sensor already knows.} Detectors fail, and the sensor ships a flag per band.
A flagged PRISMA band still occupies a slice of the cube, with $\mathrm{FWHM} = 0.0$
(\cref{ch:sensors}); it is present but meaningless.

\textbf{Stage 2, the atmosphere is opaque in places.} The five windows of \cref{ch:sensors} remove
every native band whose centre wavelength lies inside one of them, ends included.

\textbf{Stage 3, detector edges are unreliable.} A detector's response falls off at the ends of its
array, so the first and last \code{edge_bands_to_trim} bands of each detector (VNIR and SWIR) are
dropped. In the source the edge set is computed once, from the \emph{full} list of each detector's
band indices, before any other stage: a band that stage 1 or 2 already dropped can still count as
one of the three edge bands.

\textbf{Stage 4, a mostly invalid band is not worth keeping.} The per-band score $s_c$ of
\cref{eq:product:bandscore} must be at least $20\%$: a band valid in fewer than a fifth of the
scene's pixels is dropped.

\begin{worked}[label={we:bandpipeline:filter}]{Filtering a ten-band toy sensor}
A toy sensor has ten bands (all values illustrative\tagA). Bands $0$--$4$ are on detector VNIR,
$5$--$9$ on SWIR. Use \code{edge_bands_to_trim} $= 1$ to keep it small, and the real windows and
$20\%$ threshold.
\begin{center}\small
\begin{tabular}{@{}lcccccccccc@{}}\toprule
band $k$ & 0 & 1 & 2 & 3 & 4 & 5 & 6 & 7 & 8 & 9\\\midrule
wavelength (nm) & 440 & 600 & 800 & 940 & 990 & 1000 & 1400 & 1600 & 2200 & 2480\\
flag & 1 & 1 & 0 & 1 & 1 & 1 & 1 & 1 & 1 & 1\\
detector & V & V & V & V & V & S & S & S & S & S\\
$s_k$ (\%) & 99 & 99 & 99 & 99 & 99 & 99 & 99 & 12 & 99 & 99\\\bottomrule
\end{tabular}
\end{center}
First the edge set (stage 3), from the full per-detector index lists: VNIR is $\{0,\dots,4\}$, so
its edges are $0$ and $4$; SWIR is $\{5,\dots,9\}$, edges $5$ and $9$. Edge set $= \{0, 4, 5, 9\}$.
Now walk the bands.
\begin{itemize}
\item Band 0: flag 1; $440$ lies in $(0, 450)$, dropped at stage 2.
\item Band 1: passes stages 1 and 2; not an edge; $99 \ge 20$. \textbf{Kept.}
\item Band 2: flag 0, dropped at stage 1.
\item Band 3: $940$ lies in $(912, 978)$, dropped at stage 2.
\item Band 4: $990$ is in no window; but $4$ is an edge, dropped at stage 3.
\item Band 5: edge, dropped at stage 3.
\item Band 6: $1{,}400$ lies in $(1350, 1450)$, dropped at stage 2.
\item Band 7: passes 1--3; $12 < 20$, dropped at stage 4.
\item Band 8: $2{,}200$ is in no window, not an edge, $99 \ge 20$. \textbf{Kept.}
\item Band 9: edge, dropped at stage 3.
\end{itemize}
Two bands survive, $1$ and $8$. Note band $0$: it was both a window band and an edge band; it
is counted as dropped at stage 2 because the tests run in order, but because the edge set came from
the full list, trimming ``one band at the VNIR edge'' removed nothing beyond what stage 2 had already
removed. The edges are positions in the file, not ``the first surviving bands''.
\end{worked}

For a real PRISMA cube about $188$ of the $239$ bands survive stages 1--4
(\file{docs/03-data-pipeline.md}).

\section{Stages 5 and 6: deciding which pixels to keep}\label{sec:bandpipeline:pixels}

\textbf{Stage 5} applies only to EnMAP, the one sensor that ships quality masks. Any pixel flagged
($>0$) by one of the masks named in \code{quality_masks_to_apply} (default cloud, cloud shadow,
haze) has its whole validity column set to $0$: \code{validity[:, flagged] = 0}.

\textbf{Stage 6} makes validity binary. Let $C'$ be the number of bands that survived stages 1--4
and $n(i,j)$ the number of those bands in which pixel $(i,j)$ is valid. The pixel's invalid
fraction is
\begin{equation}\label{eq:bandpipeline:frac}
f(i,j) = 1 - \frac{n(i,j)}{C'},
\end{equation}
and the pixel is invalidated in every band if $f(i,j) > 0.4$
(\code{max_invalid_voxel_fraction}); note the strict ``greater than''. Pixels that pass keep their
partial validity, and stage 7 fills their holes. After stage 7, every pixel is valid in all bands or
in none: \textbf{validity is binary}. That single fact removes a whole category of special cases from
every detector downstream: none of them ever has to handle ``valid in these $140$ bands but not
those $25$''.

\begin{worked}[label={we:bandpipeline:frac}]{Kept or wiped?}
Take $C' = 188$ surviving bands. How many invalid bands may a pixel have and still be kept? It is
kept when $f \le 0.4$, that is when the number of invalid bands $u = C' - n$ satisfies
\begin{align*}
\frac{u}{188} &\le 0.4\\
u &\le 0.4 \times 188 = 75.2.
\end{align*}
$u$ is a whole number, so $u \le 75$. Check both sides of the line:
\begin{align*}
u = 75:&\quad 75/188 \approx 0.3989 \le 0.4 \;\Rightarrow\; \text{kept, its 75 holes filled at stage 7},\\
u = 76:&\quad 76/188 \approx 0.4043 > 0.4 \;\Rightarrow\; \text{invalid in every band}.
\end{align*}
\end{worked}

\begin{sourcenote}[: the denominator is the surviving bands, not 165]
The anomaly-detection course (parts 5 and 7) works this rule as ``$55/165$'' and ``$61/165$''. In the
source the fraction is taken over the bands that survived stages 1--4 (\code{num_bands =
bsq_validity.shape[0]} after band filtering, \file{prisma_dataset_builder.py}, lines 358--362;
\code{validity.mean(axis=0)} on the filtered cube in \file{common_grid_pipeline.py}). Stage 6 runs
\emph{before} resampling, when the cube has $C'$ native bands (about $188$ for PRISMA), not $165$.
The arithmetic of the course examples is right; the denominator is not.
\end{sourcenote}

\section{Linear interpolation, from the equation of a line}\label{sec:bandpipeline:linear}

Stages 7 and 8 both estimate a spectrum's value at a wavelength where it was not measured. The
simplest estimator, and the one stage 8 uses, is linear interpolation. You can derive it from the
equation of a straight line.

\subsection{Two-point form}

You know the spectrum at two wavelengths: value $y_0$ at $\lambda_0$ and $y_1$ at $\lambda_1$, with
$\lambda_0 < \lambda_1$. Assume the spectrum is a straight line between them. A straight line has a
constant \textbf{slope}\index{slope}, the rise divided by the run:
\begin{equation}\label{eq:bandpipeline:slope}
m = \frac{y_1 - y_0}{\lambda_1 - \lambda_0}.
\end{equation}
Starting at $(\lambda_0, y_0)$ and moving a distance $\lambda - \lambda_0$ to the right, the line
rises by slope times distance:
\begin{equation}\label{eq:bandpipeline:twopoint}
y(\lambda) = y_0 + \frac{y_1 - y_0}{\lambda_1 - \lambda_0}\,(\lambda - \lambda_0).
\end{equation}
This is the \textbf{two-point form}\index{linear interpolation!two-point form} of the line. Check
it at both ends. At $\lambda = \lambda_0$ the bracket is $0$, so $y = y_0$. At
$\lambda = \lambda_1$ the fraction times $(\lambda_1 - \lambda_0)$ is $y_1 - y_0$, so
$y = y_0 + y_1 - y_0 = y_1$.

\subsection{Weight form}

Name the fraction of the way from $\lambda_0$ to $\lambda_1$:
\begin{equation}\label{eq:bandpipeline:t}
t = \frac{\lambda - \lambda_0}{\lambda_1 - \lambda_0}.
\end{equation}
$t = 0$ at $\lambda_0$, $t = 1$ at $\lambda_1$, $t = 0.5$ halfway. Substitute $t$ into
\cref{eq:bandpipeline:twopoint}:
\begin{align}
y &= y_0 + (y_1 - y_0)\,t.
\intertext{Multiply out the bracket:}
y &= y_0 + y_1 t - y_0 t.
\intertext{Collect the two $y_0$ terms:}
y &= (1 - t)\,y_0 + t\,y_1. \label{eq:bandpipeline:weight}
\end{align}
This is the \textbf{weight form}\index{linear interpolation!weight form}: the estimate is a weighted
average of the two neighbours, with weights $1-t$ and $t$ that add to $1$. The nearer neighbour
gets the larger weight. It is exactly the line the product's resampler computes
(\file{spectral_resampler.py}, lines 115--118):
\begin{lstlisting}
t = np.float32((twl - wl_left) / (wl_right - wl_left))           # scalar weight t
out_2d[i, valid_cols] = (valid_spectra[idx_left] * (1.0 - t)
                         + valid_spectra[idx_right] * t)          # (n_valid,): (1-t) y0 + t y1
\end{lstlisting}
Here \code{twl} is the target wavelength and \code{valid_spectra} is \code{(C_src, n_valid)}: one
column per valid pixel, so one line computes the interpolated value for every pixel at once.

\begin{worked}[label={we:bandpipeline:linear}]{Three bands by hand}
A spectrum is known at three wavelengths: $500$, $600$ and $700$ nm, with values $1$, $2$ and $4$
(the numbers used by the repository's own interpolation test, quoted in the anomaly-detection
course, part 22). Estimate it at $550$, $650$ and $680$ nm. Do every line on paper.

\emph{At 550 nm} the neighbours are $(500, 1)$ and $(600, 2)$:
\begin{align*}
t &= (550 - 500)/(600 - 500) = 50/100 = 0.5,\\
y &= (1 - 0.5) \times 1 + 0.5 \times 2 = 0.5 + 1.0 = 1.5.
\end{align*}
\emph{At 650 nm} the neighbours are $(600, 2)$ and $(700, 4)$:
\begin{align*}
t &= (650 - 600)/(700 - 600) = 50/100 = 0.5,\\
y &= 0.5 \times 2 + 0.5 \times 4 = 1.0 + 2.0 = 3.0.
\end{align*}
\emph{At 680 nm}, same neighbours:
\begin{align*}
t &= (680 - 600)/100 = 0.8,\\
y &= 0.2 \times 2 + 0.8 \times 4 = 0.4 + 3.2 = 3.6.
\end{align*}
Check $680$ with the two-point form instead: slope $= (4 - 2)/(700 - 600) = 0.02$ per nm, and
$y = 2 + 0.02 \times 80 = 2 + 1.6 = 3.6$. Same answer, as the algebra promised.
\end{worked}

\subsection{Finding the neighbours, and what happens at the ends}

For each target wavelength the resampler finds the bracketing pair with a binary search,
\code{np.searchsorted(src_wl_sorted, twl, side="right")}, which returns the index of the first
source wavelength strictly greater than \code{twl}; call it $r$, and $\ell = r - 1$. Both are clamped
into the array: $r \le C_{\mathrm{src}} - 1$ and $\ell \ge 0$. If the clamping makes $\ell = r$,
the target lies beyond an end of the sensor's range and the resampler copies the nearest end value
(\textbf{constant extrapolation}\index{constant extrapolation}).

\begin{worked}[label={we:bandpipeline:resample}]{Resampling one spectrum onto four grid wavelengths}
A pixel's native samples (illustrative\tagA, roughly EnMAP-like spacing) are
\begin{center}
\begin{tabular}{@{}lcccc@{}}\toprule
native wavelength (nm) & 1003 & 1011 & 1019 & 1027\\
reflectance & 0.240 & 0.248 & 0.252 & 0.250\\\bottomrule
\end{tabular}
\end{center}
Resample onto the grid wavelengths $1{,}000$, $1{,}010$, $1{,}020$ and $1{,}030$ nm, following the
code's steps. Indices of the native samples are $0, 1, 2, 3$.

\emph{1,000 nm.} No native wavelength is $\le 1{,}000$; the first one greater is index $0$, so
$r = 0$, $\ell = \max(-1, 0) = 0$. $\ell = r$: constant extrapolation, $y = 0.240$.

\emph{1,010 nm.} The first native wavelength greater than $1{,}010$ is $1{,}011$ (index $1$), so
$r = 1$, $\ell = 0$:
\begin{align*}
t &= (1{,}010 - 1{,}003)/(1{,}011 - 1{,}003) = 7/8 = 0.875,\\
y &= 0.125 \times 0.240 + 0.875 \times 0.248 = 0.0300 + 0.2170 = 0.2470.
\end{align*}

\emph{1,020 nm.} First greater is $1{,}027$ (index $3$), so $r = 3$, $\ell = 2$:
\begin{align*}
t &= (1{,}020 - 1{,}019)/(1{,}027 - 1{,}019) = 1/8 = 0.125,\\
y &= 0.875 \times 0.252 + 0.125 \times 0.250 = 0.2205 + 0.03125 = 0.25175.
\end{align*}

\emph{1,030 nm.} Nothing is greater, so \code{searchsorted} returns $4$; clamped to $r = 3$, and
$\ell = 2$. Now $\ell \ne r$, so the code does \emph{not} extrapolate as a constant: it extends the
last segment's line,
\begin{align*}
t &= (1{,}030 - 1{,}019)/(1{,}027 - 1{,}019) = 11/8 = 1.375,\\
y &= (1 - 1.375) \times 0.252 + 1.375 \times 0.250 = -0.0945 + 0.34375 = 0.24925.
\end{align*}
With $t > 1$ the ``weight'' $1 - t$ is negative: this is \emph{linear} extrapolation past the last
sample, not the constant extrapolation the docstring promises. At the low end, by contrast, the
clamping produces a true constant. \Cref{fig:bandpipeline:resample} draws the four results.
\end{worked}

\begin{sourcenote}[: the resampler's upper end extrapolates linearly]
The docstring of \code{resample_to_common_grid} says it uses ``constant edge extrapolation''. Tracing
the code (\file{spectral_resampler.py}, lines 103--118) shows constant extrapolation only below the
first native wavelength; above the last, \code{idx_right} is clamped to the last index and
\code{idx_left} to the one before it, so the last segment's line is extended (with $t > 1$), as in
\cref{we:bandpipeline:resample}. On the real grid this rarely matters: the grid ends at $2{,}450$ nm
and the sensors' last bands lie close to or beyond that. It is recorded here because a reader trusting
the docstring would compute a different number.
\end{sourcenote}

\begin{figure}[htbp]\centering
\begin{tikzpicture}
\begin{axis}[width=0.9\linewidth,height=6cm,xmin=996,xmax=1034,ymin=0.236,ymax=0.256,
  xlabel={wavelength (nm)},ylabel={reflectance},axis lines=left,
  tick label style={font=\small},label style={font=\small},
  xtick={1000,1003,1010,1011,1019,1020,1027,1030},
  xticklabel style={font=\scriptsize,rotate=60,anchor=east},
  yticklabel style={/pgf/number format/fixed,/pgf/number format/precision=3},
  legend style={font=\scriptsize,at={(0.02,0.98)},anchor=north west}]
\addplot[spec,thick,mark=*,mark size=2.2pt] coordinates {(1003,0.240) (1011,0.248) (1019,0.252) (1027,0.250)};
\addlegendentry{native samples, joined by straight lines}
\addplot[only marks,mark=square*,mark size=2.4pt,loss] coordinates {(1000,0.240) (1010,0.2470) (1020,0.25175) (1030,0.24925)};
\addlegendentry{values on the common grid}
\addplot[maskgrey,dashed,forget plot] coordinates {(996,0.240) (1003,0.240)};
\addplot[loss,dashed,forget plot] coordinates {(1027,0.250) (1030,0.24925)};
\foreach \x in {1000,1010,1020,1030}{\edef\tmp{\noexpand\draw[loss!40,densely dotted] (axis cs:\x,0.236) -- (axis cs:\x,0.256);}\tmp}
\node[font=\scriptsize,anchor=north] at (axis cs:1000,0.2395) {constant};
\node[font=\scriptsize,anchor=north] at (axis cs:1029,0.2488) {linear past the end};
\end{axis}
\end{tikzpicture}
\caption{Resampling one spectrum onto the grid (\cref{we:bandpipeline:resample}). Native samples
(teal) at $1{,}003$, $1{,}011$, $1{,}019$, $1{,}027$ nm; grid values (amber) at $1{,}000$, $1{,}010$,
$1{,}020$, $1{,}030$ nm, each read off the straight line between its two neighbours. Below the first
sample the value is held constant; above the last, the code extends the last line.}
\label{fig:bandpipeline:resample}
\end{figure}

\section{Why the wavelengths must go up: a preview of a real bug}\label{sec:bandpipeline:ascending}

Binary search, and therefore both \code{np.searchsorted} and NumPy's \code{np.interp}, assumes the
sample wavelengths are in \emph{ascending} order. NumPy's documentation for \code{np.interp} says the
sequence ``is expected to be increasing, but this is not explicitly enforced'', and that otherwise
``results are meaningless'' (quoted in the anomaly-detection course, part 22). Nothing raises an
error.

PRISMA lists its bands going \emph{down} in wavelength within each detector, and its builder
concatenates SWIR before VNIR (\file{prisma_dataset_builder.py}, ``We will arrange as SWIR followd
by VNIR''). What happens when descending samples meet \code{np.interp}?

\begin{worked}[label={we:bandpipeline:descending}]{The same three samples, listed backwards}
Give \code{np.interp} the samples of \cref{we:bandpipeline:linear} in descending order:
wavelengths $(700, 600, 500)$, values $(4, 2, 1)$. Ask for $520$, $550$, $650$ and $690$ nm.

The routine's first check, made for a sorted list, is ``is the query at or beyond the last sample
wavelength?'' The last listed wavelength is $500$. Every query, $520$, $550$, $650$ and $690$, is
at least $500$, so for every query the answer is ``yes, beyond the end'', and it returns the last
listed value, $1$. All four answers are $1$
(\file{research/textbook/tools/compute_bandpipeline_output.txt})\tagC, where the correct answers are
$1.2$, $1.5$, $3.0$ and $3.8$. No single answer looks absurd; only the fact that they are all equal
gives it away. On a whole cube, every resampled band becomes a copy of one band.
\end{worked}

The product's resampler and gap-filler both sort first (\code{np.argsort} in
\file{spectral_resampler.py}, lines 70--78, and \file{spectral_interpolator.py}), so the product
was never affected. The evaluation code was: from 2026-10-17 to 2026-10-18 the evaluator's own
resample did not sort, turned every PRISMA cube it resampled into copies of one or two bands, and
rewrote a scoreboard before the bug was found and fixed. \Cref{ch:sortbug} tells that story.

\begin{keypoint}
Before any interpolation, sort the sample wavelengths ascending and reorder the values with them;
NumPy will not warn you if you forget.
\end{keypoint}

\begin{sourcenote}[: family labels are not re-ordered with the wavelengths]
\code{resample_to_common_grid} sorts the source wavelengths but passes the \emph{unsorted} family
list to \code{_assign_families} together with the sorted wavelengths
(\file{spectral_resampler.py}, lines 71 and 125). Inside, the check
\code{np.all(np.diff(source_wl_sorted) > 0)} is always true for sorted wavelengths, so the family
list is used in its original order while being indexed by sorted position. For a sensor whose
bands are not already ascending (PRISMA: SWIR first, each detector descending), the labels on the
grid come out wrong. Running the function on sources $(2000, 1500, 800, 500)$ nm with families
$(\mathrm{SWIR}, \mathrm{SWIR}, \mathrm{VNIR}, \mathrm{VNIR})$ labels $500$ nm as SWIR and $2{,}000$
nm as VNIR (\file{research/textbook/tools/compute_bandpipeline_output.txt})\tagC. Only the
bookkeeping field \code{spectral_family_order} is affected, not the reflectance values; the one
consumer found while writing this book is the destriper option of the classical
\code{StatisticalEnsembler} (\file{app/detectors/statistical_ensembler.py}, line 140), whose
behaviour with wrong labels was not traced.
\end{sourcenote}

\section{Stage 7: filling holes without inventing features}\label{sec:bandpipeline:pchip}

After stage 6 a kept pixel may still be invalid in up to $40\%$ of its bands. Stage 7 fills those
holes from the same pixel's valid bands, along wavelength. Why not a straight line here too, or a
smoother curve? The danger is a curve that \emph{overshoots}.

A \textbf{cubic spline}\index{cubic spline} joins the samples with cubic polynomials chosen to be as
smooth as possible. Smoothness has a price: given samples that rise and then level off, a spline can
swing below the lowest sample or above the highest before settling. An invented dip looks exactly
like a narrow absorption feature, which is the very evidence used to name a material.

\begin{worked}[label={we:bandpipeline:overshoot}]{A spline overshoots a step}
Four samples: $(0, 0)$, $(1, 0)$, $(2, 1)$, $(3, 1)$, a step up. With four points, SciPy's default
spline (the ``not-a-knot'' condition) is the single cubic through all four. Find it by hand. The
cubic is $0$ at $x = 0$ and $x = 1$, so it has factors $x$ and $(x - 1)$:
\[ p(x) = x(x-1)(\alpha x + \beta). \]
Use the other two points:
\begin{align*}
p(2) = 2 \times 1 \times (2\alpha + \beta) = 1 \;&\Rightarrow\; 2\alpha + \beta = 0.5,\\
p(3) = 3 \times 2 \times (3\alpha + \beta) = 1 \;&\Rightarrow\; 3\alpha + \beta = 1/6.
\end{align*}
Subtract the first equation from the second: $\alpha = 1/6 - 0.5 = -1/3$. Then
$\beta = 0.5 - 2\alpha = 0.5 + 2/3 = 7/6$. Evaluate between the samples:
\begin{align*}
p(0.5) &= 0.5 \times (-0.5) \times \left(-\tfrac{1}{3}\times 0.5 + \tfrac{7}{6}\right)
        = -0.25 \times \left(-\tfrac{1}{6} + \tfrac{7}{6}\right) = -0.25 \times 1 = -0.25,\\
p(2.5) &= 2.5 \times 1.5 \times \left(-\tfrac{1}{3}\times 2.5 + \tfrac{7}{6}\right)
        = 3.75 \times \left(-\tfrac{5}{6} + \tfrac{7}{6}\right) = 3.75 \times \tfrac{1}{3} = 1.25.
\end{align*}
The spline dips to $-0.25$, below every sample, and peaks at $1.25$, above every sample. SciPy gives
the same two numbers (\file{research/textbook/tools/compute_bandpipeline_output.txt})\tagC.
\end{worked}

\textbf{PCHIP}\index{PCHIP} (piecewise cubic Hermite interpolating polynomial) also uses one cubic
per interval, but chooses the slope at each sample to preserve the data's shape. Where the data is
flat on one side or changes direction, PCHIP sets the slope there to $0$, so it never creates a
minimum or maximum that is not in the data. On one interval $[x_0, x_1]$ with values $y_0, y_1$ and
end slopes both $0$, the Hermite cubic reduces to
\begin{equation}\label{eq:bandpipeline:hermite}
y(t) = y_0 + (y_1 - y_0)\,(3t^2 - 2t^3), \qquad t = \frac{x - x_0}{x_1 - x_0},
\end{equation}
which starts at $y_0$ ($t = 0$), ends at $y_1$ ($3 - 2 = 1$ at $t = 1$) and is flat at both ends.

\begin{worked}[label={we:bandpipeline:pchip}]{PCHIP on the same step}
The slopes between neighbouring samples (the secants) are $0$ on $[0,1]$, $1$ on $[1,2]$, $0$ on
$[2,3]$. At $x = 1$ one neighbouring secant is $0$, so PCHIP's slope there is $0$; the same at
$x = 2$; at the two ends the secants are $0$, so the end slopes are $0$. Every interval therefore uses
\cref{eq:bandpipeline:hermite}.
\begin{itemize}
\item On $[0,1]$: $y_0 = y_1 = 0$, so $y = 0$ everywhere, including $x = 0.5$. No dip.
\item On $[1,2]$ at $x = 1.5$: $t = 0.5$, $3t^2 - 2t^3 = 0.75 - 0.25 = 0.5$, so
$y = 0 + 1 \times 0.5 = 0.5$.
\item On $[2,3]$: $y_0 = y_1 = 1$, so $y = 1$, including $x = 2.5$. No bump.
\end{itemize}
SciPy's \code{PchipInterpolator} gives $0$, $0.5$ and $1$ at the three points\tagC. PCHIP gives up
some smoothness (its second derivative jumps at the samples) to guarantee it never invents a
feature. In a system that hunts for unusual spectral features, that trade is not close.
\end{worked}

The filler works per pixel on the wavelength-sorted bands, only on pixels with at least $2$ valid
bands and at least one hole; fully valid and fully invalid pixels are left alone. Holes beyond the
first or last valid band are filled with that band's value (constant extrapolation).

\begin{sourcenote}[: PCHIP only for common hole patterns]
\code{interpolate_spectral_gaps} groups pixels by their exact pattern of invalid bands and runs
PCHIP once per group. Its docstrings say groups of at least $100$ pixels get PCHIP. The code uses
\code{MIN_GROUP_SIZE = 5000} (\file{spectral_interpolator.py}, line 64); pixels whose hole pattern
is shared by fewer than $5{,}000$ pixels are filled by per-pixel \emph{linear} interpolation
(\code{np.interp}) instead. So ``stage 7 is PCHIP'' is true for common hole patterns (a dead
detector column across the scene) and false for rare ones.
\end{sourcenote}

\section{Stage 8: onto the grid}\label{sec:bandpipeline:grid}

With every kept pixel complete, stage 8 resamples each one onto the $165$ grid wavelengths by the
linear interpolation of \cref{sec:bandpipeline:linear}. The resampler's docstring gives its reason for
not using PCHIP here: ``the native bands are densely sampled (6-12nm apart) and the target grid is
10nm [...] the sample points barely move, so PCHIP offers no practical quality gain but costs
significantly more memory''. Invalid pixels stay $0$ with validity $0$ in every band; valid pixels
get validity $1$ in every band. FWHM has no meaning on the resampled grid, so
\code{band_fwhm_order} is left empty; spectral matching reads the native onboarding vendable for
that reason.

After stage 8 a PRISMA cube and an EnMAP cube have identical shape and identical wavelengths. One
model can train on shards mixing both, and a laboratory spectrum resampled onto the grid compares
against either. Resampling loses a little native detail; the pipeline accepts that deliberately,
because cross-sensor comparability is worth more.

\section{Per-band normalisation inside the model}\label{sec:bandpipeline:zscore}

The vendable's cube is called \code{normalized_hyperspectral_cube}, but that ``normalized'' meant
``cleaned by the intake pipeline''. A second, different normalisation happens inside the model, at
run time, in \file{app/foundation_models/components/pixel_normalization.py}.

\subsection{Why normalise}

Neural networks learn badly when inputs are tiny, huge, or on very different scales from each other.
Reflectance is small (around $0.04$ in the blue) and the bands differ: vegetation is dark in blue and
bright in the near infrared, so band $0$ and band $40$ have different typical values and different
spreads (\cref{fig:light:meanspec}). Fed raw, a network would spend its early training learning that
``band $40$ is usually bigger than band $0$'', a fact we already know. Normalisation tells it.

\begin{definition}[Per-band z-score]\index{z-score}\index{normalisation!per-band}
For band $c$ with mean $\mu_c$ and standard deviation $\sigma_c$ over a reference corpus, the
\textbf{z-score} of value $x_c(i,j)$ is
\begin{equation}\label{eq:bandpipeline:z}
z_c(i,j) = \frac{x_c(i,j) - \mu_c}{\sigma_c}.
\end{equation}
It says how many standard deviations the value is above ($z > 0$) or below ($z < 0$) its band's
typical value. Each band uses its own $\mu_c$ and $\sigma_c$.
\end{definition}

\begin{analogy}
Exam marks from different subjects are not comparable: a $62$ in a hard paper may be better than an
$80$ in an easy one. Converting each mark to ``standard deviations above that subject's class
average'' makes them comparable. Here each band is a subject, the corpus is the class, and the
z-score is the comparable mark. The analogy also explains the inverse: if you know a student's
z-score and the subject's average and spread, you can recover the raw mark exactly.
\end{analogy}

\begin{worked}[label={we:bandpipeline:zscore}]{Normalising the 3-band, 2-by-2 cube}
Use the cube of \cref{we:light:cube}:
\[
\mat{X}[0] = \begin{pmatrix} 0.03 & 0.05 \\ 0.05 & 0.03 \end{pmatrix},\quad
\mat{X}[1] = \begin{pmatrix} 0.20 & 0.30 \\ 0.30 & 0.20 \end{pmatrix},\quad
\mat{X}[2] = \begin{pmatrix} 0.14 & 0.10 \\ 0.14 & 0.10 \end{pmatrix}.
\]
For this exercise, take the ``corpus'' to be these four pixels and divide by $n = 4$ in the standard
deviation. Work each band separately, on paper.

\emph{Band 0.} Mean:
\begin{align*}
0.03 + 0.05 + 0.05 + 0.03 &= 0.16, \qquad \mu_0 = 0.16/4 = 0.04.
\end{align*}
Deviations $-0.01, +0.01, +0.01, -0.01$; squares $0.0001$ each; sum $0.0004$; divided by $4$,
$0.0001$; root $\sigma_0 = 0.01$.

\emph{Band 1.} $0.20 + 0.30 + 0.30 + 0.20 = 1.00$, $\mu_1 = 0.25$. Deviations $\pm 0.05$; squares
$0.0025$ each; sum $0.0100$; divided by $4$, $0.0025$; $\sigma_1 = 0.05$.

\emph{Band 2.} $0.14 + 0.10 + 0.14 + 0.10 = 0.48$, $\mu_2 = 0.12$. Deviations $\pm 0.02$; squares
$0.0004$ each; sum $0.0016$; divided by $4$, $0.0004$; $\sigma_2 = 0.02$.

\emph{Z-scores.} Pixel $(0,0)$, spectrum $(0.03, 0.20, 0.14)$:
\begin{align*}
z_0 &= (0.03 - 0.04)/0.01 = -0.01/0.01 = -1,\\
z_1 &= (0.20 - 0.25)/0.05 = -0.05/0.05 = -1,\\
z_2 &= (0.14 - 0.12)/0.02 = +0.02/0.02 = +1.
\end{align*}
The other pixels the same way:
\[
\mat{Z}[0] = \begin{pmatrix} -1 & +1 \\ +1 & -1 \end{pmatrix},\quad
\mat{Z}[1] = \begin{pmatrix} -1 & +1 \\ +1 & -1 \end{pmatrix},\quad
\mat{Z}[2] = \begin{pmatrix} +1 & -1 \\ +1 & -1 \end{pmatrix}.
\]
Raw, band $1$'s values ($0.20$ to $0.30$) were five times band $0$'s ($0.03$ to $0.05$); after
normalisation every band is centred on $0$ with spread $1$. Note what normalisation does \emph{not}
do: pixel $(0,0)$'s raw spectrum rose from band 0 to band 1, and its z-scores do not; the z-scores
describe each value relative to its own band, not the spectrum's shape.
\end{worked}

\begin{sourcenote}[: $n$ or $n-1$]
The script that produced \file{app/constants/hyperspectral.json},
\file{scripts/compute_hyperspectral_pixel_stats.py}, finishes with \code{std = np.sqrt(m2 / (n -
1))} (line 110): the \emph{sample} standard deviation, dividing by $n - 1$. With $15{,}264{,}937$
valid pixels per band\tagM{app/constants/hyperspectral.json} the difference from dividing by $n$
is about one part in thirty million, invisible. In the four-pixel example it would matter
(dividing by $3$ gives $\sigma_0 \approx 0.01155$), which is why the example states its convention.
\end{sourcenote}

\subsection{Denormalisation is the inverse map}

At the end of the model, \code{PixelDenormalize} turns the reconstruction back into reflectance with
$x = z\sigma_c + \mu_c$. Derive it by solving \cref{eq:bandpipeline:z} for $x$:
\begin{align*}
z &= \frac{x - \mu_c}{\sigma_c}.
\intertext{Multiply both sides by $\sigma_c$ (allowed because $\sigma_c > 0$):}
\sigma_c z &= x - \mu_c.
\intertext{Add $\mu_c$ to both sides:}
x &= \sigma_c z + \mu_c.
\end{align*}
Call the two maps $N(x) = (x - \mu_c)/\sigma_c$ and $D(z) = \sigma_c z + \mu_c$. Composing them in
either order gives back what you started with:
\begin{align*}
D(N(x)) &= \sigma_c \cdot \frac{x - \mu_c}{\sigma_c} + \mu_c = (x - \mu_c) + \mu_c = x,\\
N(D(z)) &= \frac{(\sigma_c z + \mu_c) - \mu_c}{\sigma_c} = \frac{\sigma_c z}{\sigma_c} = z.
\end{align*}
So $D$ is the \textbf{inverse}\index{inverse map} of $N$. The requirement $\sigma_c > 0$ is real:
a band with zero spread cannot be divided by, which is why the per-scene override below clamps tiny
spreads to $1$. With pixel $(0,0)$ band $2$ of \cref{we:bandpipeline:zscore}:
$D(+1) = 0.02 \times 1 + 0.12 = 0.14$, the original value.

The model is trained and scored in reflectance, after $D$, not in z-scores: an L1 error of $0.008$
then means ``the average band is off by $0.008$ reflectance'', a statement about the world
(\cref{ch:losses}).

\subsection{The code, and broadcasting}

\begin{lstlisting}
class PixelNormalize(nn.Module):
    def __init__(self, mean: list[float], std: list[float]):
        super().__init__()
        # (1, C, 1, 1) for broadcasting against (B, C, H, W)
        self.register_buffer("mean", torch.tensor(mean, dtype=torch.float32).view(1, -1, 1, 1))
        self.register_buffer("std",  torch.tensor(std,  dtype=torch.float32).view(1, -1, 1, 1))
    def forward(self, x):                       # x: (B, 165, H, W)
        return (x - self.mean) / self.std       # (B, 165, H, W)
\end{lstlisting}
\code{.view(1, -1, 1, 1)} reshapes the list of $165$ numbers to shape \shape{(1, 165, 1, 1)}
without changing any value (the $-1$ means ``work this size out''). Then
\textbf{broadcasting}\index{broadcasting} lines the shapes up axis by axis: where one array has size
$1$ and the other a larger size, the size-$1$ array is repeated along that axis.
\begin{center}
\begin{tabular}{@{}rl@{}}
\shape{(B, 165, H, W)} & the batch of cubes\\
$-$\ \shape{(1, 165, 1, 1)} & the means, repeated over $B$, $H$, $W$\\\midrule
\shape{(B, 165, H, W)} & each band minus its own mean
\end{tabular}
\end{center}
Band $7$'s mean is subtracted from every pixel of band $7$ in every patch, and from nothing else,
with no loop over bands.

The statistics are registered as \textbf{buffers}\index{buffer (PyTorch)}, which matters. A plain
attribute is neither saved nor trained; an \code{nn.Parameter} is saved and trained; a buffer is
saved in the checkpoint and moved to the GPU with the model, but never trained. So the $165$ means
and $165$ standard deviations travel \emph{inside} the \code{.pt} file. Normalise and denormalise
are separate modules with their own copies, so a checkpoint holds
\[ 2 \text{ modules} \times 2 \text{ statistics} \times 165 = 660 \text{ numbers}, \]
and you cannot swap the statistics without rebuilding the model: loading checks that every saved
entry has a slot of the right shape, and a length mismatch fails loudly at load time.
\Cref{ch:params} uses the $660$ to reconcile two parameter counts.

\subsection{The order of operations, and where invalid pixels land}

The model's \code{forward} masks first and normalises second
(\file{app/foundation_models/components/hyperspectral_seg_former_mae.py}, lines 126--133):
\begin{lstlisting}
if mask is not None:
    x = x * mask            # (B, 165, H, W): invalid pixels become 0.0
if self.normalize is not None:
    x = self.normalize(x)   # (B, 165, H, W): z-scores
\end{lstlisting}
An invalid pixel therefore enters the network at $z_c = (0 - \mu_c)/\sigma_c = -\mu_c/\sigma_c$, not
at $0$.

\begin{worked}[label={we:bandpipeline:real}]{Real statistics, band 0}
\file{app/constants/hyperspectral.json} gives band $0$ ($460$ nm) $\mu_0 = 0.04301$ and
$\sigma_0 = 0.03524$\tagM{app/constants/hyperspectral.json}.
\begin{align*}
\text{typical pixel, } 0.04301:&\quad (0.04301 - 0.04301)/0.03524 = 0/0.03524 = 0.00,\\
\text{bright pixel, } 0.07825:&\quad (0.07825 - 0.04301)/0.03524 = 0.03524/0.03524 = 1.00,\\
\text{dark pixel, } 0.00777:&\quad (0.00777 - 0.04301)/0.03524 = -0.03524/0.03524 = -1.00,\\
\text{invalid pixel, } 0:&\quad (0 - 0.04301)/0.03524 = -0.04301/0.03524 \approx -1.22.
\end{align*}
Over all $165$ bands, an invalid zero lands between $-2.90$ (band $27$, $730$ nm, where the mean is
large relative to the spread) and $-1.03$ (band $164$, $2{,}450$ nm)\tagC.
\end{worked}

Had the code normalised first and masked second, invalid pixels would sit at $z = 0$, which means
``perfectly typical'', indistinguishable from an ordinary pixel. Masking first puts them at a
consistent value outside the typical range, which the network can learn to recognise as ``not a
measurement''.

\subsection{When the corpus statistics do not fit: the per-scene override}

Some sensors do not deliver the units a model was fitted on. HotSat-1 ships raw 14-bit numbers
around $5{,}000 \pm 400$\tagM{app/models/training/inference_config.py}. A model whose buffers expect
values near $0.04$ would compute $z \approx (5{,}000 - 0.043)/0.035 \approx 143{,}000$: input like
nothing it has seen. \code{PixelStatsOverride} (\file{app/models/training/inference_config.py}) lets
the caller replace the buffers with the scene's own per-band mean and spread, so the input arrives
roughly centred with unit spread. The price is that scores become \emph{scene-relative}: comparable
within that scene but not across scenes, and the scoring Action records
\code{normalization_mode="per_scene_dn_zscore"} so the interface can say so. For Indradhanu the
override does not apply: band-filtered reflectance matches the training distribution, and the mode
stays \code{"baked"}.

Two details of the override are worth copying into your own code. Spreads below $10^{-6}$ (or not
finite) are clamped to $1.0$, so a constant scene cannot cause a division by zero. And the moments
are computed in 64-bit floating point, for a reason worth seeing in numbers.

\begin{worked}[label={we:bandpipeline:cancel}]{Why the moments are computed in float64}
Variance can be computed as $\E[x^2] - (\E[x])^2$, the mean of the squares minus the square of the
mean. For values $5{,}000 \pm 400$: $\E[x] = 5{,}000$, $(\E[x])^2 = 25{,}000{,}000$, and
$\E[x^2] = 25{,}000{,}000 + 400^2 = 25{,}160{,}000$. The answer, $160{,}000$, is the small difference
of two large, nearly equal numbers.

A 32-bit float carries $24$ binary digits of precision. Between $2^{24} = 16{,}777{,}216$ and
$2^{25}$ the representable numbers are $2$ apart, so a single squared value near $25{,}000{,}000$ is
stored to the nearest even integer: harmless. But the code sums squares over a million pixels. That
running sum reaches about $10^{6} \times 2.5 \times 10^{7} = 2.5 \times 10^{13}$, which lies between
$2^{44} \approx 1.76 \times 10^{13}$ and $2^{45}$, where float32 numbers are
$2^{44 - 23} = 2^{21} = 2{,}097{,}152$ apart. Each new term of about $25{,}000{,}000$ is added to a
total that can only move in steps of about two million, so each addition can be off by up to about a
million. After a million additions the sum of squares is wrong by far more than the
$160{,}000 \times 10^{6}$ that separates $\E[x^2]$ from $(\E[x])^2$ in total. The subtraction then
returns noise. This is \textbf{catastrophic cancellation}\index{catastrophic cancellation}. A 64-bit
float carries $53$ binary digits; at $2.5 \times 10^{13}$ its spacing is $2^{44 - 52} = 2^{-8}
\approx 0.004$, and the problem disappears.
\end{worked}

\begin{sourcenote}[: does the override ever fire?]
The scoring code builds the override when the vendable's \code{units} string starts with
\code{"DN_"}. As \cref{ch:product} showed, \code{VendableThermalDataset} declares no \code{units}
field, so the HotSat-1 builder's \code{units="DN_14bit_relative"} is discarded by Pydantic and
\code{getattr(vendable, "units", None)} returns \code{None} (checked by constructing one; see
\cref{sec:product:thermal}). Through this path the override therefore does not fire, despite the
comments in \file{resolver.py} and \file{_anomaly_scoring_run.py} that describe it as active.
\end{sourcenote}

\begin{drift}
\emph{Dated 2026-10-29.} Indradhanu V2 reads native bands and normalises with
\code{WavelengthStats} (same file, \file{pixel_normalization.py}): the mean and spread are stored as
curves over wavelength, measured on SpectralEarth, and read at each sensor's own band wavelengths by
linear interpolation, after sorting the curve because SpectralEarth's band order is not monotonic
near $900$--$1{,}000$ nm. It is \cref{eq:bandpipeline:z} with $\mu$ and $\sigma$ looked up by
wavelength instead of by position (\cref{ch:crossstem,ch:pretrain}).
\end{drift}

\coursefig[0.9\linewidth]{iw-09-normalise}{Why each band is normalised separately (course figure,
Indradhanu part 9). (a) Raw reflectance puts a dim blue band and a bright near-infrared band on very
different scales. (b) Subtracting each band's mean and dividing by its standard deviation puts every
band on one scale, centred at $0$ with spread $1$. Distributions are illustrative.}{fig:bandpipeline:normalise}

\section{The nearest-valid pixel fill}\label{sec:bandpipeline:fill}

After the eight stages, invalid pixels hold $0$ in every band. The \code{band_filter_apply} Action
then (by default, \code{apply_nearest_valid_fill: true}) replaces each invalid pixel's values with
those of its nearest valid pixel, leaving the validity cube unchanged
(\file{backend/allotrope/action_types/band_filter_apply.py}, lines 372--381;
\file{app/utils/pixel_fill/nearest_valid_fill.py}).

\begin{definition}[Nearest-valid fill]\index{nearest-valid fill}
For each invalid pixel $(i,j)$, find the valid pixel $(i^*, j^*)$ that minimises the Euclidean
distance $\sqrt{(i - i^*)^2 + (j - j^*)^2}$, and copy all $C$ band values from $(i^*, j^*)$ to
$(i,j)$. Valid pixels are their own nearest valid pixel (distance $0$) and are unchanged. Ties are
broken by the implementation, SciPy's \code{distance_transform_edt}.
\end{definition}

The motivation given in the module's docstring: a patch embedding whose kernel straddles a boundary
between valid pixels and invalid zeros sees a cliff, produces contaminated tokens there, and
attention spreads the contamination, so boundaries light up in the error map and outrank real
detections. Copying a real neighbouring measurement turns the cliff into a smooth continuation, and
every filled value is physically plausible because it is a genuine measurement, never an average or
an interpolation. The masks are untouched, so any loss still excludes those pixels.

\begin{worked}[label={we:bandpipeline:fill}]{Filling a row and a 3-by-3 patch}
\emph{One row} (the docstring's example; one band): values $30, 32, 0, 0, 0, 25, 28$ with validity
$1, 1, 0, 0, 0, 1, 1$. Positions $2$, $3$, $4$ are invalid.
\begin{itemize}
\item Position 2: valid neighbours at distance $1$ (position $1$, value $32$) and $3$ (position $5$).
Nearest is position 1: $32$.
\item Position 3: distance $2$ to position $1$ and $2$ to position $5$, a tie. The implementation
picked position $1$: $32$.
\item Position 4: distance $1$ to position $5$: $25$.
\end{itemize}
Result $30, 32, 32, 32, 25, 25, 28$\tagC.

\emph{A 3-by-3 patch} (one band, illustrative\tagA):
\[
\text{values} = \begin{pmatrix} 0.10 & 0.12 & 0 \\ 0.11 & 0 & 0 \\ 0 & 0 & 0.30 \end{pmatrix},\qquad
\text{validity} = \begin{pmatrix} 1 & 1 & 0 \\ 1 & 0 & 0 \\ 0 & 0 & 1 \end{pmatrix}.
\]
\begin{itemize}
\item $(0,2)$: nearest valid is $(0,1)$ at distance $1$: $0.12$.
\item $(1,1)$: $(0,1)$ and $(1,0)$ both at distance $1$, a tie; the implementation picked $(1,0)$:
$0.11$.
\item $(1,2)$: $(2,2)$ at distance $1$ ($(0,2)$ is itself invalid): $0.30$.
\item $(2,0)$: $(1,0)$ at distance $1$: $0.11$.
\item $(2,1)$: $(2,2)$ at distance $1$: $0.30$.
\end{itemize}
\[
\text{filled} = \begin{pmatrix} 0.10 & 0.12 & 0.12 \\ 0.11 & 0.11 & 0.30 \\ 0.11 & 0.30 & 0.30 \end{pmatrix}
\]
(confirmed by running the function, \file{research/textbook/tools/compute_bandpipeline_output.txt})\tagC.
\end{worked}

\begin{sourcenote}[: the kernel, the cliff, and whether the fill reaches Indradhanu]
Three corrections, all from reading the source.
\begin{enumerate}
\item \emph{The kernel.} The fill's docstring, the anomaly-detection course (part 7) and
\file{docs/03-data-pipeline.md} speak of the SegFormer's ``7x7 kernel''. The encoder that both
Chakshu and Indradhanu use builds its first patch embedding with kernel $4$ and stride $4$, and the
later stages with kernel $3$ (\file{app/foundation_models/components/seg_former_encoder.py}, lines
110--111). The $7$ is the default in \file{overlap_patch_embedding.py}'s docstring, not what the
encoder passes. \Cref{ch:patchembed} covers the patch embedding.
\item \emph{The cliff.} ``Invalid pixels normalise to about $-2.3$'' is the thermal model's number
(an invalid $0\,^\circ$C pixel). For Indradhanu an invalid zero lands at $-1.22$ in band $0$ and
between $-2.90$ and $-1.03$ across bands (\cref{we:bandpipeline:real}).
\item \emph{The re-zeroing.} The scoring code passes the model \code{mask_tensor} built from the same
spatial validity used by the fill (\file{_anomaly_scoring_run.py}, line 286), and the model's
\code{forward} multiplies by that mask before normalising (lines 127--128 of
\file{hyperspectral_seg_former_mae.py}; the inferencer calls \code{self.model(tensor, mask=mask, ...)},
\file{segformer_mae_inferencer.py}, line 222). The filled pixels are therefore set back to $0$
before Indradhanu v0.2.0 sees them, and enter at $-\mu_c/\sigma_c$ as in
\cref{we:bandpipeline:real}. As far as this book could trace, the fill does not change Indradhanu
v0.2.0's input on this path; it does change what reaches code that reads the cube without
re-masking it (the evaluator also applies it, \file{app/evaluation/datasets.py}, line 136). Which
consumers benefit was not established here.
\end{enumerate}
\end{sourcenote}

\begin{confusions}
\confusion{Does resampling lose information?}{Yes, a little, deliberately: cross-sensor
comparability is worth more than the last few nanometres of native resolution. Spectral matching
avoids the loss by reading the native onboarding vendable.}
\confusion{Why not interpolate across the atmospheric windows too?}{There is no signal there to
interpolate from; the gap is up to $170$ nm wide. Filling it would invent spectrum. The grid keeps the
hole.}
\confusion{188 bands survive stage 4 but the grid has 165. Where did the rest go?}{Nowhere: they are
different axes. The $188$ are surviving sensor bands at sensor wavelengths; the $165$ are grid slots.
Stage 8 resamples one onto the other.}
\confusion{Is the 0.4 in stage 6 a fraction of the cube or of the pixel?}{Of that pixel's surviving
bands. A pixel with more than $40\%$ of its own bands invalid is wiped.}
\confusion{Is PCHIP or linear interpolation used?}{Both. Stage 7 (gap filling) uses PCHIP for common
hole patterns and linear for rare ones; stage 8 (onto the grid) is linear.}
\confusion{Is the vendable's ``normalized'' cube the z-score?}{No. The vendable holds reflectance; the
z-score is computed inside the model from statistics stored in the checkpoint.}
\confusion{Why denormalise at the end instead of comparing in z-score space?}{So the loss and the
score are in reflectance units, which mean something physically.}
\end{confusions}

\begin{checkyourself}
\item\label{ex:bandpipeline:1} Linear interpolation: a spectrum is $0.20$ at $1{,}600$ nm and $0.26$
at $1{,}612$ nm. Estimate it at $1{,}610$ nm using the weight form, then check with the two-point
form. Show every step.
\item\label{ex:bandpipeline:2} With $C' = 200$ surviving bands and \code{max_invalid_voxel_fraction}
$= 0.4$, what is the largest number of invalid bands a kept pixel can have? Is a pixel with exactly
$80$ invalid bands kept? Why does the answer depend on ``$>$'' versus ``$\ge$''?
\item\label{ex:bandpipeline:3} Give \code{np.interp} wavelengths $(900, 800, 700)$ and values
$(0.3, 0.2, 0.1)$, and ask for $750$ nm. What does it return, and what is the correct answer? Explain
the wrong one using the ``beyond the last sample'' check.
\item\label{ex:bandpipeline:4} A band has corpus mean $0.25$ and standard deviation $0.05$. Normalise
the values $0.15$, $0.25$, $0.40$ and an invalid pixel stored as $0$. Then denormalise $z = -0.6$ and
show algebraically that denormalising the z-score of $0.40$ returns $0.40$.
\item\label{ex:bandpipeline:5} Apply nearest-valid fill (one band) to the row
$0, 0.20, 0, 0, 0, 0.35$ with validity $0, 1, 0, 0, 0, 1$. Which position is a tie? Then explain in two sentences why the fill
copies a measurement rather than averaging the neighbours.
\end{checkyourself}

\nextchapter{ch:linalg}
